% !TeX root = ../../Book3.tex
\section{正规族}
\label{b3:sec:0402}

研究一个函数时，积分公式恢复它的内部值；研究一族函数时，同一公式控制整个族的变化。正规性要求任意序列都有可用的子列，因而是一种紧性性质，而不是每个成员的光滑程度。

\subsection{定义}
\label{b3:subsec:040201}

\begin{definition}\label{b3:def:04-normal-finite}
设 \(\mathcal F\) 是区域 \(U\) 上的全纯函数族。若任取 \(f_n\in\mathcal F\)，都能抽取在 \(U\) 上局部一致收敛至有限全纯函数的子列，则称 \(\mathcal F\) 在有限值意义下是\term[zhenggui-zu]{正规族}[normal family]。
\end{definition}
极限不必属于原族，所以这里对应相对序列紧性。族 \(\{z^n:n\ge1\}\) 在圆盘上正规，其某些子列极限为0，但0原本不在该族内。族 \(\{nz:n\ge1\}\) 在包含0的圆盘上不正规：若局部一致有有限极限，导数在0也应收敛，而相应导数为无界整数。

\ProseParagraphBreak

有些教材允许子列趋于无穷。为了处理极点，下面改用球面距离统一表述；有限值版本和球面版本不能未经说明互换。

\subsection{序列紧性}
\label{b3:subsec:040202}

局部一致拓扑可以由穷尽紧集定义度量
\[
 \mathrm{d}_{\mathrm{loc}}(f,g)=\sum_{j=1}^{\infty}2^{-j}\min\{1,\sup_{K_j}|f-g|\}.
\]
级数始终收敛；三角不等式从每个截断上确界得到。该度量趋于0当且仅当所有 \(K_j\) 上一致趋于0：前向由第 \(j\) 项的下界得到，反向先使尾和小，再控制有限个前项。因此“抽子列”与这个度量空间的序列紧性相对应。

\ProseParagraphBreak

对角线抽取的顺序不能省略。先在 \(K_1\) 选收敛子列，再在其内选 \(K_2\) 上的子列，依次进行，最后取第 \(j\) 列的第 \(j\) 个成员。对任何固定 \(K_m\)，对角线尾部都属于第 \(m\) 次子列，所以仍一致收敛；各紧集上的极限在重叠处一致，拼成 \(U\) 上同一个函数。若只对每点分别抽子列，就不能保证得到了同一条子列。

\ProseParagraphBreak

局部正规性能够拼成全域正规性。
\begin{proposition}\label{b3:prop:04-normal-locality}
设 \(\mathcal F\) 为开集 \(U\) 上的全纯函数族。若每点都有邻域，使限制族在该邻域有限值正规，则整个族在 \(U\) 上有限值正规。
\end{proposition}
\begin{proof}
平面第二可数性使这些邻域有可数子覆盖 \(V_1,V_2,\ldots\)。对任意函数序列，先在 \(V_1\) 抽取局部一致子列，再依次在 \(V_j\) 抽取，最后用对角线得到每个 \(V_j\) 上局部一致收敛的同一子列。重叠处的点值极限唯一，所以局部极限拼成 \(U\) 上全纯函数。若 \(K\subset U\) 紧，可用有限个位于各 \(V_j\) 内的小闭盘覆盖 \(K\)，这些盘上的一致收敛给出全 \(K\) 一致收敛。
\end{proof}
可数性与紧性承担不同作用：可数子覆盖允许一次对角线抽取，紧集的有限子覆盖把局部误差合成一个全紧集误差。只说“每点可以抽子列”而不经过这两个步骤，还没有得到正规性。

\subsection{球面度量}
\label{b3:subsec:040203}

扩充平面的\term[qiumian-juli]{球面距离}[chordal distance]取为
\[
 \chi(z,w)=\frac{|z-w|}{\sqrt{1+|z|^2}\sqrt{1+|w|^2}},\qquad
 \chi(z,\infty)=\frac1{\sqrt{1+|z|^2}}.
\]
它是单位球面弦长的一半，经立体投影得到，故满足三角不等式。归一化常数不会改变正规性，但应固定一种写法。倒数变换满足 \(\chi(1/z,1/w)=\chi(z,w)\)，包括0和无穷远。

\ProseParagraphBreak

在有限值有界集合上，球面距离与欧氏距离等价：若 \(|z|,|w|\le M\)，则 \(|z-w|/(1+M^2)\le\chi(z,w)\le|z-w|\)。趋向无穷则表示模一致变大。函数 \(n\) 的常值列没有有限局部极限，却在球面意义下一致趋向 \(\infty\)。这一点使两种正规性明确不同。

\subsection{亚纯函数族}
\label{b3:subsec:040204}

\begin{definition}\label{b3:def:04-meromorphic-normal}
区域 \(U\) 上的亚纯函数族若每个序列都有球面局部一致收敛子列，其极限是亚纯函数或常值 \(\infty\)，则称它在球面意义下正规。
\end{definition}
\begin{proposition}\label{b3:prop:04-spherical-limit}
亚纯函数列的球面局部一致极限，若存在，则是亚纯函数或常值 \(\infty\)。
\end{proposition}
\begin{proof}
球面一致极限先是到球面的连续映射 \(f\)。若 \(f(a)\) 有限，缩小圆盘使像避开 \(\infty\) 的一个邻域；充分大的 \(f_n\) 在这个圆盘也有限且一致有界，球面收敛转成欧氏收敛，Weierstrass 定理给出 \(f\) 全纯。若 \(f(a)=\infty\)，在小圆盘中像避开0，充分大的 \(1/f_n\) 全纯有界并欧氏一致收敛到 \(1/f\)。其极限全纯。若极限不恒零，零点孤立，对倒数取回得到 \(f\) 亚纯；若它在一个圆盘恒零，\(f\) 在该盘恒为无穷。全局若同时有有限值点，则连接两种区域的边界附近，倒数极限的恒等定理迫使继续为零，产生矛盾。连通性因此使恒无穷情形传播到整个区域。
\end{proof}
族 \(nz\) 在0附近没有球面一致极限，因为点0处始终为0，任意非零点却趋于无穷，候选极限不连续。允许无穷极限仍不能消除局部尖峰。

\ProseParagraphBreak

球面一致收敛允许极点的位置移动。例子 \(f_n(z)=1/(z-1/n)\) 在每个有限平面紧集上球面一致趋于 \(1/z\)：对有限非极点，
\[
 \chi\left(\frac1{z-1/n},\frac1z\right)
 =\chi(z-1/n,z)\le\frac1n,
\]
等式利用倒数是球面等距变换，并且在移动极点处也按连续性成立。欧氏一致收敛则不可能在含0的紧集上成立，因为函数有极点。球面坐标正是为了让这种自然极限连续可见。

\ProseParagraphBreak

全纯函数列的球面极限若不是恒为无穷，则实际上仍全纯，不会凭空产生一个孤立极点。若极限在 \(a\) 为极点，取小盘使其倒数在盘内全纯、仅中心有零点。充分晚的 \(f_n\) 在小盘内远离0，所以 \(1/f_n\) 全纯无零，局部一致趋于有中心零点但不恒零的 \(1/f\)，与 Hurwitz 无零极限结论矛盾。因此极点只能来自亚纯成员或整域无穷退化。这个细节在不同正规族定义之间换用时尤其重要。
